\subsection{开集的截面；闭集的截面；开集的投影。部分连续性}

命题 4. 设 \(\mathbf{X}_1, \mathbf{X}_2\) 是两个拓扑空间；则对每个 \(a_1 \in \mathbf{X}_1\)，映射 \(x_2 \to (a_1, x_2)\) 是 \(\mathbf{X}_2\) 映到 \(\{a_1\} \times \mathbf{X}_2\)（\(\mathbf{X}_1 \times \mathbf{X}_2\) 的子空间）上的一个同胚。

这是把第 1 目命题 1 的推论 1 应用于常值函数 \(x_{2} \rightarrow a_{1}\) 而得的一个特殊情形。

映射 \(x_{2} \to (a_{1}, x_{2})\) 是关于等价关系 \(\mathrm{pr}_{2}z = \mathrm{pr}_{2}z'\)（在 \(\mathbf{X}_{1} \times \mathbf{X}_{2}\) 中）的一个连续截面（§ 3，第 5 目）；因此 \(\mathbf{X}_{1} \times \mathbf{X}_{2}\) 按此等价关系的商空间同胚于 \(\mathbf{X}_{2}\)。

推论. 截面 \(\mathbf{A}(x_1)\)（开（相应地，闭）集 \(\mathbf{A}\)——积 \(\mathbf{X}_1 \times \mathbf{X}_2\) 的——在任意一点 \(x_1 \in \mathbf{X}_1\)（*）处的截面）在 \(\mathbf{X}_2\) 中是开（相应地，闭）集。

命题 5. 积 \(X_{1} \times X_{2}\) 的开集 U 到任一因子上的投影是开集。

例如，我们有 \(\mathrm{pr}_2\mathrm{U} = \bigcup_{x_1\in X_1}\mathrm{U}(x_1)\)，于是由命题 4 的推论及公理 \((\mathbf{O}_{\mathbf{I}})\) 即得该命题。

\pandocbounded{\includegraphics[keepaspectratio]{images/part001_fc188d82.jpg}}

注 1. \(\mathbf{X}_1 \times \mathbf{X}_2\) 的闭子集到一个因子上的投影未必是闭集。例如，在有理平面 \(\mathbf{Q}^2\) 中，方程为 \(x_1 x_2 = 1\) 的双曲线是闭集，但它的两个投影都等于 \(\mathbf{Q}\) 中点 0 的余集，而这不是闭集。

命题 6. 设 \(\mathbf{X}_1, \mathbf{X}_2, \mathbf{Y}\) 是三个拓扑空间，\(f\) 是积空间 \(\mathbf{X}_1 \times \mathbf{X}_2\) 到 \(\mathbf{Y}\) 中的一个映射。若 \(f\) 在点 \((a_1, a_2)\) 处连续，则部分映射 \(x_2 \to f(a_1, x_2)\)（\(\mathbf{X}_2\) 到 \(\mathbf{Y}\) 中的）在点 \(a_2\) 处连续。

因为这个映射是 \(f\) 与映射 \(x_{2} \to (a_{1}, x_{2})\) 的复合，故结论由命题 4 得出。

命题 6 常被表述为：两个变量的连续函数关于每个变量分别连续。

\pandocbounded{\includegraphics[keepaspectratio]{images/part001_31ca3096.jpg}}

注 2. 由映射 \(f: \mathbf{X}_{1} \times \mathbf{X}_{2} \to \mathbf{Y}\) 所确定的所有部分映射都连续，而 \(f\) 在 \(\mathbf{X}_{1} \times \mathbf{X}_{2}\) 上不连续，这是可能的（参见第九章，§ 5，习题 23）。* 例如，若 \(f\) 是实平面 \(\mathbb{R}^{2}\) 到 \(\mathbb{R}\) 中的映射，其定义为

\[
f (x, y) = x y / \left(x ^ {2} + y ^ {2}\right) \quad \text { if } \quad (x, y) \neq (0, 0)
\]

且 \(f(0, 0) = 0\)，则所有部分映射都连续；但 \(f\) 在 \((0, 0)\) 处不连续，因为 \(f(x, x) = 1/2\)（对 \(x \neq 0\)）。

若 \(g\) 是 \(\mathbf{X}_1\) 到 \(\mathbf{Y}\) 中的一个映射，在点 \(a_1\) 处连续，则映射 \((x_1, x_2) \to g(x_1)\)（\(\mathbf{X}_1 \times \mathbf{X}_1 \to \mathbf{Y}\)）在所有点 \((a_1, x_2)\) 处连续，因为它是 \(g\) 与到 \(\mathbf{X}_1\) 上的投影的复合。

本小节的结果容易推广到拓扑空间的任意积 \(\prod_{i\in I}X_i\)，只需注意到对 I 的任意分拆 (J, K)，这个积同胚于积 \(\left(\prod_{i\in J}X_i\right)\times\left(\prod_{i\in K}X_i\right)\)（第 1 目，命题 2）。

